\documentclass[11pt]{article}
\usepackage{mathunal-base}
\mathunalfuente{serif}

\renewcommand{\examtitulo}{Primer Examen Parcial de Álgebra Lineal}
\renewcommand{\examfecha}{10 de septiembre de 2022}

\newcommand{\vf}{\fbox{V}\ \fbox{F}}
\newcommand{\sino}{\fbox{Sí}\ \fbox{No}}

\begin{document}

\examheader

\vspace{4pt}
\noindent
\begin{minipage}[t]{0.56\linewidth}
\textbf{Puntaje. Sólo para uso Oficial}\\[4pt]
\begin{tabular}{|c|c|c|c|c|c|c|}
\hline
1-3 & 4-5 & 6 & 7 & 8 & \textbf{TOTAL} & \textbf{NOTA} \\ \hline
 & & & & & & \\ \hline
\end{tabular}
\end{minipage}%
\hfill
\begin{minipage}[t]{0.4\linewidth}
Nombre: \dotfill\\[6pt]
Cédula: \dotfill\\[6pt]
Profesor: \dotfill\\[6pt]
Grupo: \dotfill
\end{minipage}
\vspace{6pt}\par
\noindent\rule{\linewidth}{0.4pt}

\vspace{6pt}
\noindent\textbf{Instrucciones:} La duración del examen es de 1 hora y 50 minutos. El examen consta de ocho preguntas en tres hojas impresas por ambos lados, verifique que su examen esté completo. En las preguntas con procedimiento justifique sus respuestas en los espacios asignados. No está permitido sacar hojas en blanco ni ningún tipo de apuntes durante el examen, verifique que su celular esté apagado. No se permite el uso de calculadora.

\vspace{6pt}
\noindent\textbf{IDENTIFICACIÓN}

\vspace{10pt}
\noindent\textbf{I. Completación} En las preguntas 1 a 6 complete los espacios en blanco. \textbf{NOTA}: En esta sección se califica sólo la respuesta y no se tiene en cuenta el procedimiento.

\begin{enumerate}
\item{}[10pt] Sean $A = \begin{bmatrix}7 & 4\\ 1 & -1\end{bmatrix}$ y $b = \begin{bmatrix}1\\-1\end{bmatrix}$. Encuentre $||Ab||^2$.

\vspace{4pt}
\fbox{\begin{minipage}[t][2.5cm]{\dimexpr\linewidth-2\fboxsep-2\fboxrule}
$||Ab||^2 = $
\end{minipage}}

\item{}[10pt] Sean $B = \begin{bmatrix}1 & 0 & 1\\ 1 & 1 & 0\end{bmatrix}$ y $C = \begin{bmatrix}2 & 2\\ 2 & 2\\ 2 & 2\end{bmatrix}$. Calcule $\mathrm{Rango}(BC)$.

\vspace{4pt}
\fbox{\begin{minipage}[t][2.5cm]{\dimexpr\linewidth-2\fboxsep-2\fboxrule}
$\mathrm{Rango}(BC) = $
\end{minipage}}

\item{}[12pt] Determine si las siguientes proposiciones son verdaderas o falsas (Marque V o F según el caso).
\begin{enumerate}
\item{}[3pt] Todo sistema de ecuaciones no homogéneo con más variables que ecuaciones es consistente.

\vf

\item{}[3pt] Todo sistema homogéneo con más ecuaciones que variables es inconsistente.

\vf

\item{}[3pt] Sea $V=\{v_1,v_2,v_3\}$ un conjunto con tres vectores de $\mathbb{R}^3$. Entonces $gen(V)=\mathbb{R}^3$.

\vf

\item{}[3pt] Si una matriz cuadrada $A$ es simétrica entonces $A^T$ es simétrica.

\vf
\end{enumerate}

\item{}[10pt] Considere la matriz $D = \begin{bmatrix}a & b & c\\ d & e & f\\ g & h & i\end{bmatrix}$ y suponga que $\det(2D^T) = -24$. Calcule los siguientes determinantes.

\vspace{6pt}
\noindent (a) [5pt]

\vspace{4pt}
\fbox{\begin{minipage}[t][3cm]{\dimexpr\linewidth-2\fboxsep-2\fboxrule}
$\begin{vmatrix} -g & -h & -i\\ 2d+5a & 2e+5b & 2f+5c\\ a & b & c\end{vmatrix} = $
\end{minipage}}

\vspace{8pt}
\noindent (b) [5pt]

\vspace{4pt}
\fbox{\begin{minipage}[t][2cm]{\dimexpr\linewidth-2\fboxsep-2\fboxrule}
$\det(3D^2) = $
\end{minipage}}

\item{}[13pt] Sea $A = \begin{bmatrix}1 & 0 & 3\\ k & 1 & 6\\ 1 & 2 & 3\end{bmatrix}$
\begin{enumerate}
\item{}[5pt] Encuentre todos los valores de $k$ tales que $\det(A) = -12$.

\vspace{2.5cm}

\item{}[8pt] Sea $b = \begin{bmatrix}0\\4\\8\end{bmatrix}$. ¿Existen valores de $k$ tales que el sistema $[A|b]$ sea inconsistente? Escriba sí o no según el caso.

\vspace{4pt}
\sino
\end{enumerate}

\item{}[10pt] Considere los vectores $v_1 = \begin{bmatrix}1\\1\\0\\1\end{bmatrix}$, $v_2 = \begin{bmatrix}1\\0\\1\\0\end{bmatrix}$ y $v_3 = \begin{bmatrix}2\\1\\1\\1\end{bmatrix}$ y sea $A = [v_1|v_2|v_3]$ la matriz de tamaño $4\times 3$ cuyas columnas son los vectores $v_1,v_2,v_3$.
\begin{enumerate}
\item{}[5pt] ¿Son los vectores $v_1, v_2$ y $v_3$ linealmente independientes?

\vspace{4pt}
\sino

\item{}[5pt] ¿Tiene el sistema homogéneo asociado a $A$ una única solución?

\vspace{4pt}
\sino
\end{enumerate}
\end{enumerate}

\newpage
\noindent\textbf{\large II. Solución con Procedimiento}

\vspace{6pt}
\noindent En esta parte del examen se deben justificar las respuestas.

\begin{enumerate}
\setcounter{enumi}{6}
\item{}[10pt] Sean $u, v$ vectores en $\mathbb{R}^n$. Suponga que el ángulo entre $u$ y $v$ es $\frac{\pi}{3}$ y que $||u||=1$ y $||v||=1$. Calcule $||u-v||$.

\vspace{4cm}

\item{}[25pt] Considere la siguiente red de transporte:

\vspace{6pt}
\begin{center}
\begin{tikzpicture}[>=Stealth, scale=1]
  \coordinate (A) at (0,2);
  \coordinate (B) at (0,-2);
  \coordinate (C) at (3.6,0);
  \draw[thick] (0,3.2) -- (0,-3.2);
  \draw[thick,->] (A) -- node[above,sloped]{$f_2$} (C);
  \draw[thick,->] (A) -- node[left]{$f_1$} (B);
  \draw[thick,->] (B) -- node[below,sloped]{$f_3$} (C);
  \fill (A) circle (2pt) node[above right]{A};
  \fill (B) circle (2pt) node[below right]{B};
  \fill (C) circle (2pt) node[right]{C};
  % flujos externos
  \draw[->] (0,3.2) -- (0,2.35) node[midway,right]{12};
  \draw[->] (-1.3,3.1) -- (A) node[midway,above left]{8};
  \draw[->] (-1.3,-1.3) -- (B) node[midway,above left]{4};
  \draw[->] (B) -- (0,-3.2) node[midway,right]{16};
  \draw[->] (C) -- (4.9,1.1) node[midway,above]{4};
  \draw[->] (C) -- (4.9,-1.1) node[midway,below]{4};
\end{tikzpicture}
\end{center}

\begin{enumerate}
\item{}[12pt] Establezca y resuelva un sistema de ecuaciones lineales para encontrar los flujos posibles en la red de la figura anterior.

\vspace{5cm}

\item{}[7pt] Si las direcciones de los flujos se respetan, ¿cuáles son los posibles flujos mínimos y máximos a través de cada rama?

\vspace{4cm}

\item{}[6pt] Suponiendo que las direcciones de los flujos se respetan, halle los valores de $f_1, f_2$ y $f_3$ de tal forma que se obtenga el mínimo flujo vehicular en la rama $\overline{AB}$.
\end{enumerate}
\end{enumerate}

\examfooter
\end{document}
